\begin{document}
\maketitle

\begin{abstract}
This paper isolates the analytic closure layer of TFPT. The selector $\Padm$ is treated as a
physical admissible-sector construction, while the dynamics is carried by $Z_{\mathrm{rel}}$,
admissible Schwinger distributions, Osterwalder--Schrader reconstruction, the local Minkowski net,
stable massive scattering, and the exact admissible RG flow.
\end{abstract}

\begin{tfptfrontbox}
\textbf{Inputs from previous papers.} Paper 1 supplies $\Pprim$. Paper 2 supplies the carrier and
discrete determinant data needed for $P_{\mathrm{sing}}$ and $P_\Theta$. Paper 3 is not logically
required except for optional observable cross-references.

\textbf{New theorem contribution.} Conditional nonperturbative closure on the admissible sector:
\[
\Padm=\Pprim P_{\mathrm{sing}}P_\Theta,\qquad
\theta_{\mathrm{eff}}=0,\qquad
\operatorname{arg}\det M_u=\operatorname{arg}\det M_d=0,
\]
together with reflection positivity, OS reconstruction, local Minkowski net, stable massive
scattering, and exact admissible flow.

\textbf{Not claimed here.} No $\alpha$ detail calculation, no CMB, no E8, and no full empirical
tables.

\textbf{Falsification or audit surface.} The paper fails if selector and dynamics are conflated, if
positivity/gap hypotheses are hidden, or if strong-CP closure uses an inadmissible phase convention.
\end{tfptfrontbox}

\begin{tfptwarningbox}
Preferred claim language: \emph{conditional nonperturbative closure under explicitly stated
admissibility, positivity, and gap hypotheses}. Avoid unqualified ``proved QFT'' rhetoric.
\end{tfptwarningbox}

\begin{tfptwarningbox}
Every analytic theorem in this paper is to be read with its hypotheses in the statement, not as an
unconditional construction of four-dimensional QFT. The standing checklist is: admissible
Schwinger family, temperedness, Euclidean covariance, reflection positivity, clustering,
locality-domain regularity, and a stated massive-sector gap where scattering is invoked.
\end{tfptwarningbox}


\begin{tfptclaimcontract}
\textbf{Claim.} $\Padm$ selects the physical sector; strong CP closes; QFT reconstruction is conditional on stated analytic hypotheses.\par
\textbf{Inputs.} $\Pprim$, carrier/determinant data, admissibility complex, positivity/gap hypotheses.\par
\textbf{First assumptions.} Reflection positivity, temperedness, Euclidean covariance, clustering, locality-domain regularity, sector gap.\par
\textbf{Proof status.} Conditional nonperturbative closure / standard-theorem interface.\par
\textbf{Kill condition.} Hidden positivity/gap assumption, nonzero strong-CP angle, or failure of OS/local-net hypotheses.\par
\end{tfptclaimcontract}

\begin{tfptstatusbox}
\textbf{\stag{L} (Lean verified) used as input.}\ The $3+2$
hypercharge carrier from Paper 2 supplies $Y$,
$(q_-,q_+) = (-2,3)$, $\tr_E Y = 0$, and the carrier polynomial.

\smallskip
\textbf{\stag{I} (Lean interfaced).}\ The Paper-1 boundary
input is \texttt{BoundaryPolarization}; the seam-class input
is \texttt{SeamWindingData}; the Yukawa Stage-D upstream input
is \texttt{BoundaryYukawaKernel}. The body's manuscript
"primitive Yukawa generator" / "branch Yukawa tensor"
construction is intended to factor through these typed
targets.

\smallskip
\textbf{\stag{M} (manuscript-theorem content).}\ The
manuscript-level theorems in the source-extract body
(Yukawa rigidity from the primitive component of the branch
Yukawa tensor, primitive Yukawa generator existence, lowest
primitive Yukawa generator) are \stag{M}.

\smallskip
\textbf{\stag{C} (conditional, the QFT closure layer).}\ The
admissibility selector $\Padm$, the strong-CP closure
$\theta_{\mathrm{eff}} = 0$, the OS / reflection-positivity
reconstruction, the local-net structure, the scattering /
S-matrix layer, and the exact RG / nonperturbative flow are
\emph{conditional} on the analytic hypotheses listed in the
claim contract. They are \emph{not} \stag{L} and should be
read with their hypotheses in the statement.
\end{tfptstatusbox}

\tableofcontents

\section{Selector Versus Dynamics}

The central distinction is
\[
\Padm\quad\text{selects the physical sector},
\]
whereas dynamics is carried by
\[
Z_{\mathrm{rel}}
\Rightarrow
\{S_n^T\}
\Rightarrow
(\mathcal H_{\mathrm{adm}},\mathfrak A_{\mathrm{adm}})
\Rightarrow
\Gamma_k
\Rightarrow
\GrenTFPT .
\]
This separation must be repeated in the introduction and conclusion because it is the main defense
against overclaiming.

\section{Full Admissibility Complex}

After carrier and determinant data are fixed, the full selector is
\[
\Padm=\Pprim P_{\mathrm{sing}}P_\Theta.
\]
The paper should retain the two-stage upgrade from $\Pprim$ to $\Padm$, the Hodge projector
description, and the exact role of the gap projection. Carrier material is summarized only as an
input.

\section{Strong CP Closure}

The strong-CP sector is stated as an admissibility result:
\[
\theta_{\mathrm{eff}}=0,
\qquad
\operatorname{arg}\det M_u=0,
\qquad
\operatorname{arg}\det M_d=0.
\]
The argument should connect hadronic singlet selection, determinant structure,
$\gamma_5$-Hermiticity, and the sheet involution without importing phenomenological tuning.

\section{Reflection Positivity and OS Reconstruction}

The reflection-positivity block should state all assumptions on the admissible Schwinger functions,
fermionic positivity, quotient form, and domain regularity. The OS theorem is used as a standard
interface only after the admissible-sector positivity hypotheses are named.

\section{Local Net and Scattering}

The local Minkowski net, Haag--Ruelle/LSZ massive scattering, and dressed massless interface are
treated as conditional closure outputs. The paper must distinguish stable massive scattering from
low-curvature dressed massless behavior.

\section{Exact Admissible Flow}

The exact admissible RG flow is the analytic continuation of the same sector:
\[
\partial_k\Gamma_k
=
\frac12\operatorname{STr}\left[(\Gamma_k^{(2)}+R_k)^{-1}\partial_kR_k\right]_{\mathrm{adm}},
\]
with the admissible projection included as part of the flow definition. The observable hierarchy is
included only as far as it supports the QFT closure layer.

\clearpage
\section{Main Technical Development}

This section contains the main technical development assigned to this paper by the TFPT 4.5 clean split. Cross-paper background is referenced through dependency and scope boxes; extended backend material is kept in the Technical Companion.

\begin{tfptwarningbox}
\textbf{Status-tier mapping for the source-extract body.}\
The text below carries the manuscript precursors of the
Yukawa Stage-D content. Specifically, the
``\emph{Primitive Yukawa generator}'' definition, the
``\emph{Lowest primitive Yukawa generator on the closed
branch}'' lemma, and the ``\emph{Carrier rigidity from the
primitive component of the branch Yukawa tensor}'' theorem
are the source-side picture of what the companion Lean note
records as the v11 typed interface
\[
\text{branch Yukawa tensor (manuscript)}
\;\Longrightarrow\;
\texttt{BoundaryYukawaKernel}\,K\,E_-
\;\Longrightarrow\;
\bigl(\omega,\ \texttt{CI},\ \texttt{CS},\ \texttt{Nontrivial}\,E_-\bigr).
\]
Read the body under this mapping. The
\stag{L} content of the rank-3 discharge is
\texttt{YukawaStageDExistence.stage\_D\_iff\_finrank\_eq\_three};
the manuscript theorem ``Carrier rigidity from the primitive
component of the branch Yukawa tensor'' is \stag{M}, intended
to factor through the \stag{I} target
\texttt{BoundaryYukawaKernel}. Strong-CP, hadronic admissibility,
and seam-inflow theorems are \stag{C} (conditional on the
analytic hypotheses listed in the claim contract).
\end{tfptwarningbox}

\tfptDraftXrefNotice

\input{source_extracts/04_qft_source.tex}

\section{Source Extraction Map}

\begin{tfptsourcebox}
Use \sourceDraft:
\begin{itemize}[leftmargin=1.5em]
    \item Sections 7.4 and 7.5 for hadronic admissibility and strong-CP ingredients.
    \item Section 8 as the main technical body, with carrier repetition removed.
    \item Sections 8.20--8.22 for topological positivity and strong-CP closure.
    \item Sections 8.23--8.41 for reflection positivity, OS closure, local net, and scattering.
    \item Section 13.2 for exact admissible RG flow, 1PI expansion, and observable hierarchy.
    \item Appendices E, F, and H as the technical backend.
\end{itemize}
\end{tfptsourcebox}


\begin{tfptexports}
Exports: $\Padm=\Pprim P_{\mathrm{sing}}P_\Theta$, $\theta_{\mathrm{eff}}=0$, admissible Schwinger/OS interface, local-net and RG-flow hypotheses/results under conditional closure.
\end{tfptexports}

\section{Not Used Here}

Exact electromagnetic benchmark derivations, CKM/PMNS prediction tables, gravity/metrology
readouts, CMB spectra, sky realization, E8 grammar, full empirical ledgers, horizons, and transient
channels are not used as proof inputs in this paper.

\end{document}
